\subsubsection{Inversion de Hensel dans la tour arithmétique et réalisation par histoires}
\label{sec:ch-capacidad-hensel}

Avant de relever les branches en histoires enrichies, il convient de retrouver
l’argument fini qui détermine la capacité de la tour arithmétique. On
incorpore la preuve complète du transducteur non filtré de la matrice ; son
domaine est distingué de l’arbre survivant.

\begin{proposition}[Conjugaison de la tour de Hensel non filtrée]
\label{prop:ch-hensel-bruto}
Soient \(p\) premier, \(d\geq1\) et \(A\in\operatorname{GL}_d(\mathbb Z)\). Pour des vecteurs lignes, en prenant des représentants de résidus dans \(\{0,\ldots,p-1\}^d\), on définit
\[
 x_tA=u_t+px_{t+1},\qquad u_t=x_tA\bmod p.
\]
L'application
\[
 \Psi_T:(\mathbb Z/p^T\mathbb Z)^d\longrightarrow
             (\mathbb F_p^d)^T,\qquad
 x_0\longmapsto(u_0,\ldots,u_{T-1})
\]
est bijective, d’inverse
\[
 x_0\equiv\sum_{t=0}^{T-1}p^tu_tA^{-(t+1)}\pmod{p^T}.
\]
Les bijections commutent avec les réductions et donnent un homéomorphisme
\(\Psi:\mathbb Z_p^d\to(\mathbb F_p^d)^{\mathbb N}\), où
\(\mathbb Z_p:=\varprojlim_T\mathbb Z/p^T\mathbb Z\). L’actualisation
\(x_t\mapsto x_{t+1}\) correspond à supprimer le premier bloc.
\end{proposition}

\begin{proof}
L'inverse entier de \(A\) existe car son déterminant est \(\pm1\). En itérant \(x_t=(u_t+px_{t+1})A^{-1}\), on obtient
\[
 x_0=\sum_{t=0}^{T-1}p^tu_tA^{-(t+1)}+p^Tx_TA^{-T}.
\]
Modulo \(p^T\), un mot détermine une unique classe initiale. Étant donné un
mot arbitraire, fixer \(x_T=0\) et reconstruire à rebours produit une
classe qui l’émet, ce qui prouve la surjectivité. La même formule
prouve la compatibilité lors de la troncature. Le passage à la limite inverse donne
l’homéomorphisme et la relation avec le décalage du ruban.
\end{proof}

Pour \(p=3\) et \(d=6\), l’alphabet de blocs a \(3^6=729\)
éléments. Cette capacité arithmétique n’est pas identifiée par définition à
la réalisabilité de chaque bloc après chaque préfixe survivant. Le
lien avec les arêtes \(\operatorname{Ext}\) est précisément la fonction
du lemme~\ref{lem:elevacion-catalogal-nonadica}. Les deux
opérations que distingue la matrice sont ainsi maintenues : représentation dans la tour brute et
réalisation sous les compatibilités APP--TRIT--TPK.
